% %--------------------------
% %	CONFIGURATION
% %--------------------------
\documentclass[11pt,a4paper]{moderncv}
\moderncvstyle{classic}
\moderncvcolor{black}
\definecolor{firstnamecolor}{rgb}{0,0,0}

% Basic packages
\usepackage[utf8]{inputenc}
\usepackage[T1]{fontenc}
\usepackage[hscale=0.7, top=2cm, bottom=2cm]{geometry}
\usepackage[dvipsnames]{xcolor}
\definecolor{DarkBlue}{RGB}{0, 0, 150}
\usepackage[colorlinks=true, urlcolor=DarkBlue, linkcolor=DarkBlue]{hyperref}

\usepackage{microtype}

% Configure font settings
\renewcommand{\rmdefault}{ppl}
\renewcommand{\sfdefault}{phv}
\renewcommand*{\namefont}{\fontsize{26}{18}\mdseries}

% Configure microtype
\microtypesetup{expansion=false}

% Load sensitive information
\input{../../other_material/personal_info.tex}

% %--------------------------
\begin{document}
\makecvtitle
Hello Electricity Maps team,

\vspace*{2em}
I am writing to apply for the Senior Data Scientist, Price Forecasting role in Copenhagen. I am a data scientist with seven years of experience across forecasting, production machine learning, pharmaceutical analytics, healthcare AI, and energy-related consulting work. The role interests me because it combines several things I care about: statistical modelling, real production systems, and work where better forecasts can support better decisions in a complex physical market.

\hspace*{2em}
My strongest direct fit is forecasting. At Signum Life Science, my current work is on an end-to-end Databricks forecasting platform for pharmaceutical price and demand forecasting, including patient population forecasts. I also maintain and improve a Denmark-specific XGBoost price forecasting model for fortnightly blind pharmaceutical auctions. The domain is different from electricity, but the practical problem is similar in the ways that matter to me: noisy market behaviour, limited observability, changing incentives, and the need for models that can be trusted enough to guide decisions.

\hspace*{2em}
I have also built and maintained production ML systems beyond forecasting. At Valcon, I architected and implemented an end-to-end production pipeline for a real-time neural network using natural-language inputs and high-dimensional labels, and worked on data science projects across pharma, energy, IoT, and manufacturing. I also mentored colleagues, led explainable AI and Git trainings, and regularly explained modelling work to non-technical stakeholders. Those experiences are relevant to a senior role where model quality, engineering practice, and communication all matter.

\hspace*{2em}
My background is grounded in statistics and applied mathematics. I studied Engineering Physics and Statistical Learning and Data Analytics at KTH, with coursework including Statistical Machine Learning, Applied Time Series, and Optimization, and later spent an exchange period in Applied Mathematics at ETH. Earlier, at RaySearch Laboratories, I worked on probabilistic modelling and optimisation for automated radiotherapy planning, using autoencoders and sparse Gaussian processes on segmented 3D CT scans. That work made me comfortable with models whose outputs need to be useful in real decision workflows, not only accurate on paper.

\hspace*{2em}
I do not want to overstate my fit: I have not worked directly with electricity price forecasting, BigQuery, or grid and weather datasets. What I can offer is strong adjacent experience in price forecasting, production ML, MLflow and Python-based data science, stakeholder-facing model delivery, and enough energy-domain exposure to make the transition realistic. Electricity Maps is especially interesting to me because it applies data science to electricity systems and climate impact rather than to optimisation for its own sake.

\vspace*{2em}
I would welcome the opportunity to discuss how my forecasting and production ML background could be useful to the team. Regards,

\vspace*{1em}
Ivar Cashin Eriksson.

\end{document}
